\section{Levantamento de custos}
\label{sec:custos}

Este capítulo organiza os custos de operação do \projeto{} em \textbf{fixos} e \textbf{variáveis},
com valores estimados em reais. O modelo considera a fase de lançamento hiperlocal em Belo
Horizonte (Atividade~III) e o \textbf{mês~12} de operação como referência de regime.

\subsection{Premissas e legenda}

\begin{itemize}
  \item \textbf{Cenário A (enxuto):} sócios sem pró-labore; apenas despesas operacionais em caixa.
  \item \textbf{Cenário B (equipe remunerada):} soma ao cenário A uma ajuda de custo ao gestor de
  rede e um pró-labore simbólico aos quatro sócios.
  \item \textbf{Câmbio:} R\$~\cambio{} por US\$ (a PTAX de 02/10/2026 foi R\$~5,2238).
  \item \textbf{Origem de cada número:} \bP{} pesquisa em fonte pública (ver Referências);
  \bH{} hipótese do grupo a validar; \bC{} valor a cotar ou confirmar com fornecedor ou contador.
\end{itemize}

% ---------------------------------------------------------------------
\subsection{Custos fixos}
\label{sec:custos-fixos}

Despesas que não variam diretamente com o número de usuários ativos ou atendimentos intermediados.

\begin{longtable}{L{5.5cm} R{2.1cm} C{1.6cm} L{3.8cm}}
\toprule
\cabfundo \cab{Item} & \cab{R\$/mês} & \cab{Base} & \cab{Observação} \\
\midrule
\endfirsthead
\toprule
\cabfundo \cab{Item} & \cab{R\$/mês} & \cab{Base} & \cab{Observação} \\
\midrule
\endhead
\bottomrule
\endfoot
Contabilidade e formalização (ME/Simples) & 300,00 & \bC & Cotar com contador \\
Infraestrutura de backend & 100,00 & \bP\,\bH & Firebase Blaze ou VPS para Flask \\
Google Maps Platform (busca geolocalizada) & 100,00 & \bH & Depende do uso da API de busca \\
Domínio, e-mail e ferramentas de produto & 65,00 & \bH & Figma, GitHub e Canva \\
Consultoria veterinária (curadoria da triagem) & 600,00 & \bH & Pode cair com parceria acadêmica \\
Jurídico e LGPD (R\$~3.000 em 12 meses) & 250,00 & \bC & Termos, privacidade e parecer CFMV \\
Divulgação local e mídias impressas & 150,00 & \bH & Cartazes e QR codes (Atividade~8) \\
Conta de desenvolvedor Google Play (US\$~25 em 12 meses) & 10,83 & \bP & Taxa única, diluída em 1 ano \\
\midrule
\totalfundo \textbf{Total do cenário A} & \textbf{1.575,83} & & \\
\midrule
Ajuda de custo do gestor de rede (meio período) & 1.500,00 & \bH & Exclusivo do cenário B \\
Pró-labore simbólico (4 sócios $\times$ R\$~1.000) & 4.000,00 & \bH & Exclusivo do cenário B \\
\midrule
\totalfundo \textbf{Total do cenário B} & \textbf{7.075,83} & & \\
\end{longtable}

% ---------------------------------------------------------------------
\subsection{Custos variáveis}
\label{sec:custos-variaveis}

Despesas operacionais que crescem diretamente com o volume de uso e transações na plataforma.

\begin{longtable}{L{4.4cm} R{2.8cm} C{1.6cm} L{3.9cm}}
\toprule
\cabfundo \cab{Item} & \cab{Valor} & \cab{Base} & \cab{Observação} \\
\midrule
\endfirsthead
\toprule
\cabfundo \cab{Item} & \cab{Valor} & \cab{Base} & \cab{Observação} \\
\midrule
\endhead
\bottomrule
\endfoot
Taxa de pagamento via Pix & 0,99\% & \bP & Sobre a fração paga em dinheiro \\
Taxa de pagamento via cartão à vista & 4,98\% & \bP & Sobre a fração paga em dinheiro \\
\textbf{Taxa média de gateway} (60\% Pix / 40\% cartão) & 2,59\% & \bH & Mix médio estimado para intermediação \\
Imposto sobre receita (Simples Nacional) & 6,00\% & \bC & Anexo III (sujeito ao Fator R $\ge$ 28\%; atinge 15,5\% no Anexo V sem pró-labore) \\
Taxa da loja sobre assinaturas (Google Play) & 15,00\% & \bC & Incide \textbf{apenas} sobre o PetCuida+ (consultas usam gateway externo) \\
Infraestrutura por usuário ativo & R\$~0,10/MAU & \bH & Custos de banco de dados e tráfego extra \\
Suporte e moderação por usuário ativo & R\$~0,20/MAU & \bH & Atendimento ao cliente e curadoria \\
Fundo de liquidação de créditos solidários & R\$~6,60/atend. & \bH & 40\% dos atendimentos usam crédito; repasse de 50\% do valor de face à clínica \\
\end{longtable}

O \emph{fundo de liquidação} é o custo mais sensível do modelo: é o recurso que remunera as
clínicas parceiras pelos créditos aceitos, conforme a estratégia de mitigação de risco da Atividade~III
(``liquidação parcial em dinheiro e garantia de repasse por campanhas patrocinadas'').

% ---------------------------------------------------------------------
\subsection{Visão consolidada no mês 12}
\label{sec:consolidado}

Premissas de volume ($\bH$): 1.500 usuários ativos por mês (MAU); 10\% realizam um atendimento
intermediado por mês (150 atendimentos); ticket médio de R\$~110,00 com 70\% pagos em dinheiro; 5\% assinam o
PetCuida+ (75 assinantes); 10 clínicas parceiras ativas; R\$~1.500,00/mês provenientes de convênios e patrocínio social.

\begin{table}[htbp]
\centering
\caption{Receitas, custos e resultado estimados para o mês 12 (R\$).}
\label{tab:consolidado}
\begin{tabular}{L{8.3cm} R{3.0cm}}
\toprule
\cabfundo \cab{Linha} & \cab{R\$/mês} \\
\midrule
Comissão sobre atendimentos (150 $\times$ R\$~7,70) & 1.155,00 \\
Assinaturas PetCuida+ (75 $\times$ R\$~12,90) & 967,50 \\
Plano para clínicas (10 $\times$ R\$~79,90) & 799,00 \\
Convênios e patrocínio social & 1.500,00 \\
\midrule
\totalfundo \textbf{Receita total} & \textbf{4.421,50} \\
Taxas de pagamento e da loja & (483,60) \\
Impostos sobre receita & (265,29) \\
Infraestrutura e suporte & (450,00) \\
Fundo de liquidação de créditos & (990,00) \\
\midrule
\totalfundo \textbf{Margem de contribuição} (50,5\%) & \textbf{2.232,61} \\
Custos fixos -- cenário A & (1.575,83) \\
\totalfundo \textbf{Resultado -- cenário A} (margem líquida de 14,9\%) & \textbf{656,78} \\
Custos fixos -- cenário B & (7.075,83) \\
\totalfundo \textbf{Resultado -- cenário B} & \textbf{(4.843,22)} \\
\bottomrule
\end{tabular}
\end{table}

No cenário A o negócio atinge o ponto de equilíbrio com receita modesta devido aos baixos custos fixos. No
cenário B, com equipe remunerada, a operação exige maior escala de usuários ou captação de fomento externo (ponto de equilíbrio detalhado na Seção~\ref{sec:equilibrio}).